\documentclass[a4paper,12pt]{article}
\usepackage[T2A]{fontenc}
\usepackage[utf8]{inputenc}
\usepackage[russian]{babel}
\usepackage{amsmath,amssymb,mathtools}
\usepackage{cmbright}
\renewcommand{\familydefault}{\sfdefault}
\usepackage{icomma}
\usepackage{geometry}
\geometry{left=23mm,right=23mm,top=22mm,bottom=24mm}
\usepackage{microtype}
\usepackage{xcolor}
\usepackage{tikz}
\usetikzlibrary{calc,arrows.meta,positioning,shapes.geometric}
\usepackage{tabularx,booktabs}
\usepackage{enumitem}
\usepackage{parskip}
\usepackage{fancyhdr}
\usepackage{setspace}
\usepackage{hyperref}
\hypersetup{hidelinks}

\definecolor{accent}{HTML}{176B70}
\definecolor{darkaccent}{HTML}{194F55}
\definecolor{textgray}{HTML}{303638}
\color{textgray}
\setstretch{1.10}
\setlength{\headheight}{16pt}
\pagestyle{fancy}
\fancyhf{}
\fancyhead[L]{\small Ксения Васильченко}
\fancyhead[R]{\small 20.10.26}
\fancyfoot[C]{\small\thepage}
\renewcommand{\headrulewidth}{0.35pt}
\setlist[enumerate]{leftmargin=*,itemsep=.8em,topsep=.4em}
\newcommand{\sectionline}[1]{\section*{\color{darkaccent}#1}}
\newcommand{\learn}[1]{\par\noindent$\square$\enspace #1\par}
\newcommand{\keyidea}[1]{\par\noindent\textcolor{darkaccent}{\textbf{Запомните.}}\enspace #1\par}
\newcommand{\mix}[3]{#1\,\frac{#2}{#3}}

\begin{document}
\hypersetup{pageanchor=false}
\begin{titlepage}
\thispagestyle{empty}
\centering
\vspace*{30mm}
{\color{accent}\large МАТЕМАТИКА \textbullet\ 6 КЛАСС\par}
\vspace{18mm}
{\color{darkaccent}\Huge\bfseries Чек-лист\par}
\vspace{9mm}
{\LARGE Умножение дробей\par}
\vspace{3mm}
{\Large и смешанных чисел\par}
\vspace{8mm}
{\large Вычисления со скобками \textbullet\ Текстовые задачи\par}
\vspace{19mm}
{\large 20 октября 2026 года\par}
\vfill
{\large Лёвин Артём Александрович\par}
\vspace{3mm}
{\normalsize эксклюзивно для Ксении Васильченко\par}
\vspace{34mm}
\begin{tikzpicture}[remember picture,overlay]
\draw[accent!58,line width=1.25pt]
  ([xshift=25mm,yshift=20mm]current page.south west)
  .. controls ([xshift=79mm,yshift=42mm]current page.south west)
           and ([xshift=-77mm,yshift=4mm]current page.south east)
  .. ([xshift=-25mm,yshift=20mm]current page.south east);
\end{tikzpicture}
\end{titlepage}
\hypersetup{pageanchor=true}

\sectionline{1. Умножение обыкновенных дробей}
\learn{Вы умножаете числитель на числитель, а знаменатель на знаменатель.}
\learn{Вы различаете умножение и сложение дробей: общий знаменатель при умножении не нужен.}

Для $b\ne0$ и $d\ne0$ выполняется правило:
\[
\boxed{\frac ab\cdot\frac cd=\frac{a\cdot c}{b\cdot d}}.
\]
\textbf{Пример 1.} Умножим две дроби:
\[
\frac23\cdot\frac45=\frac{2\cdot4}{3\cdot5}=\boxed{\frac8{15}}.
\]

\textbf{Пример 2.} Сократим \emph{перед} перемножением:
\[
\frac23\cdot\frac94
=\frac{2\cdot9}{3\cdot4}
=\frac{1\cdot3}{1\cdot2}
=\frac32=\boxed{\mix{1}{1}{2}}.
\]
Сокращать можно \emph{любой множитель числителя с любым множителем
знаменателя}. В примере числа $2$ и $4$ сокращаются на $2$, а $9$ и $3$ --- на $3$.

\keyidea{Если после умножения получилась неправильная дробь, выделите целую часть.
Сокращайте дробь при первой удобной возможности: так меньше вычислений.}

\sectionline{2. Дробь, натуральное число и степень}
\learn{Вы представляете любое натуральное число в виде дроби со знаменателем $1$.}
\learn{Вы возводите дробь в квадрат как произведение двух одинаковых дробей.}

\textbf{Пример 3.} Умножим дробь на целое число:
\[
\frac23\cdot9
=\frac23\cdot\frac91
=\frac{2\cdot3}{1\cdot1}=\boxed6.
\]
Любое натуральное число $n$ можно записать как $\frac n1$;
знаменатель $1$ можно затем опустить: $\frac61=6$.

\textbf{Пример 4.} Возведём дробь в квадрат:
\[
\left(\frac23\right)^2
=\frac23\cdot\frac23
=\boxed{\frac49}.
\]
Степень относится ко всей дроби, поэтому записывайте дробь в скобках.

\clearpage
\sectionline{3. Несколько множителей и действия в скобках}
\learn{Вы сокращаете множители до вычисления большого произведения.}
\learn{Вы выполняете действия в скобках в первую очередь.}

\textbf{Пример 5.} При умножении трёх дробей действует то же правило:
\[
\frac13\cdot\frac35\cdot\frac{15}{4}
=\frac{1\cdot3\cdot15}{3\cdot5\cdot4}
=\boxed{\frac34}.
\]

\textbf{Пример 6.} Сначала сложим дроби в скобках, затем умножим:
\begin{align*}
\left(\frac49+\frac5{12}\right)\cdot\frac{18}{31}
&=\left(\frac{16}{36}+\frac{15}{36}\right)\cdot\frac{18}{31}\\
&=\frac{31}{36}\cdot\frac{18}{31}
=\boxed{\frac12}.
\end{align*}
\keyidea{Для \emph{сложения и вычитания} нужен общий знаменатель;
для \emph{умножения} перемножают числители и знаменатели. В комбинированном выражении
сначала завершите действия в скобках.}

\sectionline{4. Умножение смешанных чисел}
\learn{Вы переводите смешанное число в неправильную дробь.}
\learn{Вы получаете окончательный ответ в сокращённом виде.}

\[
\boxed{\mix{a}{b}{c}=\frac{a\cdot c+b}{c}}
\qquad (a\text{ --- целая часть},\;0\le b<c).
\]

\textbf{Пример 7.} Умножим смешанное число на дробь:
\[
\mix{2}{3}{4}\cdot\frac37
=\frac{2\cdot4+3}{4}\cdot\frac37
=\frac{11\cdot3}{4\cdot7}
=\frac{33}{28}=\boxed{\mix{1}{5}{28}}.
\]

\textbf{Если есть сложение или вычитание смешанных чисел в скобках,}
сначала выполните соответствующее действие; вспомните приём «занять единицу»,
если дробной части уменьшаемого недостаточно.

\textbf{Пример 8.} Вычислим по действиям:
\begin{align*}
\left(\mix{1}{1}{2}-\frac56\right)\cdot\frac94
&=\left(\frac32-\frac56\right)\cdot\frac94\\
&=\left(\frac96-\frac56\right)\cdot\frac94
=\frac23\cdot\frac94=\boxed{\mix{1}{1}{2}}.
\end{align*}

\clearpage
\sectionline{5. Текстовая задача: умножение по смыслу}
\learn{Вы находите дробь от числа умножением.}
\learn{Вы отличаете величину «на столько-то больше» от «во столько-то раз больше».}

\keyidea{Чтобы найти $\frac mn$ от числа $A$, вычислите $A\cdot\frac mn$.
Если дана масса одного кубического сантиметра, умножьте её на объём тела.}

\textbf{Пример 9.} Объём свинцового шарика равен $\mix{3}{1}{3}\ \text{см}^3$.
Масса $1\ \text{см}^3$ свинца равна $\mix{11}{17}{50}\ \text{г}$.
Найдём массу шарика:
\[
\mix{3}{1}{3}\cdot\mix{11}{17}{50}
=\frac{10}{3}\cdot\frac{567}{50}
=\frac{189}{5}=\boxed{\mix{37}{4}{5}\ \text{г}}.
\]

\sectionline{6. Задачи на встречное движение}
\learn{Вы составляете таблицу «скорость --- время --- расстояние».}
\learn{Вы учитываете, что при движении навстречу пути \emph{складываются}.}

\[
\boxed{s=v\cdot t},\qquad
v=\frac st\ (t\ne0),\qquad
t=\frac sv\ (v\ne0).
\]

\begin{center}
\begin{tabular}{@{}lccc@{}}
\toprule
& Скорость, км/ч & Время, ч & Путь, км \\
\midrule
Поезд 1 & $75$ & $2$ & $75\cdot2=150$\\
Поезд 2 & $\frac45\cdot75=60$ & $2$ & $60\cdot2=120$\\
\bottomrule
\end{tabular}
\end{center}

\textbf{Пример 10.} Поезда вышли одновременно из двух городов навстречу друг другу.
Второй движется со скоростью $\frac45$ от скорости первого, равной $75$ км/ч.
Они встретились через $2$ часа. Расстояние между городами:
\[
S=150+120=\boxed{270\ \text{км}}.
\]
\begin{center}
\begin{tikzpicture}[>=Stealth,scale=1]
\coordinate (A) at (0,0); \coordinate (M) at (4.1,0); \coordinate (B) at (7.4,0);
\draw[very thick,accent] (A)--(B);
\foreach \p in {A,M,B}{\fill[darkaccent] (\p) circle (2.1pt);}
\node[above=4pt] at (A) {Город A};
\node[above=4pt] at (M) {Встреча};
\node[above=4pt] at (B) {Город B};
\draw[->,very thick,accent] (.45,-.48)--(3.55,-.48);
\draw[<-,very thick,accent] (4.6,-.48)--(6.96,-.48);
\node[below=15pt] at (2.05,0) {$150$ км};
\node[below=15pt] at (5.75,0) {$120$ км};
\end{tikzpicture}
\end{center}

\sectionline{7. Одинаковое направление: найдите разность путей}
\learn{Вы рисуете начало движения и конечное положение каждого объекта.}
\learn{Вы вычитаете пройденные пути, если объекты стартовали вместе и движутся в одном направлении.}

\enlargethispage{1\baselineskip}
\textbf{Мини-пример.} Жаба прыгает на $6$ см, лягушка --- на $15$ см.
После $10$ прыжков каждая пройдёт $60$ см и $150$ см соответственно.
Расстояние между ними: $150-60=\boxed{90\ \text{см}}$.
\clearpage
\thispagestyle{plain}
\begin{center}
{\color{darkaccent}\LARGE\bfseries Алгоритм: как умножать дроби и смешанные числа}
\end{center}
\vspace{8mm}
\begin{center}
\begin{tikzpicture}[
  node distance=11mm,
  >={Stealth[length=2.7mm]},
  every node/.style={font=\small,align=center},
  terminator/.style={draw=darkaccent,rectangle,rounded corners=5pt,minimum width=54mm,minimum height=11mm,fill=accent!6,thick},
  action/.style={draw=accent,rounded corners=2mm,text width=81mm,minimum height=13mm,inner sep=3mm,thick},
  choice/.style={draw=darkaccent,diamond,aspect=3.3,text width=43mm,inner xsep=1mm,inner ysep=1mm,minimum height=14mm,thick},
  path/.style={->,line width=.95pt,darkaccent}
]
\node[terminator] (start) {Начало};
\node[action,below=of start] (convert) {Смешанные числа переведите в неправильные дроби,\\ натуральные числа --- в дроби со знаменателем $1$};
\node[choice,below=13mm of convert] (reduce) {Есть общие множители\\ в числителе и знаменателе?};
\node[action,below=13mm of reduce] (multiply) {Перемножьте числители и знаменатели};
\node[action,right=5mm of reduce,text width=35mm,minimum height=14mm] (cancel) {Сократите\\ общие множители};
\node[choice,below=13mm of multiply] (improper) {Результат ---\\ неправильная дробь?};
\node[action,right=5mm of improper,text width=35mm,minimum height=14mm] (extract) {Выделите\\ целую часть};
\node[terminator,below=14mm of improper] (finish) {Ответ};
\draw[path] (start)--(convert);
\draw[path] (convert)--(reduce);
\draw[path] (reduce)--node[left,pos=.55]{нет}(multiply);
\draw[path] (reduce)--node[above]{да}(cancel);
\draw[path] (cancel.south) |- ([yshift=6mm]multiply.east) -- (multiply.east);
\draw[path] (multiply)--(improper);
\draw[path] (improper)--node[left,pos=.55]{нет}(finish);
\draw[path] (improper)--node[above]{да}(extract);
\draw[path] (extract.south) |- ([yshift=6mm]finish.east) -- (finish.east);
\end{tikzpicture}
\end{center}

\clearpage
\sectionline{Тренировочный блок --- Путь А}
\textbf{Выполните 10 заданий.} В вычислениях сначала проверяйте, можно ли
сократить множители. Окончательные дроби сокращайте, из неправильных
по возможности выделяйте целую часть. В задачах запишите решение и ответ.

\begin{enumerate}
\item Вычислите: $\displaystyle\frac23\cdot\frac45$.
\item Вычислите с сокращением: $\displaystyle\frac57\cdot\frac{14}{20}$.
\item Вычислите: $\displaystyle\frac45\cdot20$.
\item Возведите дробь в квадрат: $\displaystyle\left(\frac23\right)^2$.
\item Перемножьте три дроби: $\displaystyle\frac27\cdot\frac{14}{8}\cdot\frac54$.
\item Вычислите по действиям:
$\displaystyle\left(\frac49+\frac5{12}\right)\cdot\frac{18}{31}$.
\item Выполните действия со смешанными числами:
$\displaystyle\left(\mix{2}{3}{4}-\mix{1}{1}{4}\right)\cdot\frac37$.
\item Масса одного кубического сантиметра вещества ---
$\displaystyle\mix{11}{17}{50}$ г. Найдите массу тела объёмом
$\displaystyle\mix{3}{1}{3}$ см$^3$.
\item Два поезда выехали одновременно навстречу друг другу и встретились
через $\displaystyle\mix{2}{1}{2}$ часа. Скорость первого равна
$72$ км/ч, второго --- $\frac56$ скорости первого.
Найдите расстояние между городами. Составьте таблицу «скорость --- время --- путь».
\item Жаба и лягушка одновременно начали прыгать с одного места
в одном направлении. Длина прыжка жабы --- $\displaystyle\mix{7}{1}{5}$ см,
а прыжок лягушки в $\displaystyle\mix{2}{1}{3}$ раза длиннее.
На каком расстоянии они окажутся, если каждая сделает по $15$ прыжков?
Сделайте схематический рисунок.
\end{enumerate}

\clearpage
\sectionline{Ответы к тренировочному блоку}
\vspace{5mm}
\renewcommand{\arraystretch}{1.7}
\begin{center}
\begin{tabularx}{\linewidth}{@{}c>{\raggedright\arraybackslash}X@{}}
\toprule
\textbf{№} & \textbf{Ответ}\\
\midrule
1 & $\dfrac8{15}$\\
2 & $\dfrac12$\\
3 & $16$\\
4 & $\dfrac49$\\
5 & $\dfrac58$\\
6 & $\dfrac12$\\
7 & $\dfrac9{14}$\\
8 & $\mix{37}{4}{5}$ г\\
9 & $330$ км\\
10 & $144$ см\\
\bottomrule
\end{tabularx}
\end{center}

\vspace{8mm}
\keyidea{Проверьте вычисления обратной оценкой: при умножении на правильную
положительную дробь число уменьшается; при умножении на число больше $1$
увеличивается. В текстовой задаче убедитесь, что единицы измерения ответа верны.}

\end{document}
